\section{数据与奖励：训练质量的第一约束}

\subsection{数据结构从 ``样本'' 变为 ``轨迹''}
Agent 训练中的数据不只是问答对，而是环境状态、动作、工具调用结果、延迟反馈与最终 outcome 的组合轨迹。课程反复提醒，如果数据只保留最终答案，模型会失去 ``如何到达答案'' 的学习信号。

\begin{knowledgebox}{轨迹数据的四个必要字段}
\begin{enumerate}
    \item 上下文状态：任务背景、可用工具、历史动作。
    \item 决策动作：模型当前选择的操作与参数。
    \item 环境反馈：工具返回、执行日志、外部错误信息。
    \item 终局标签：成功/失败、质量评分、代价统计。
\end{enumerate}
\end{knowledgebox}

\subsection{奖励设计：局部奖励与终局奖励的平衡}
课程给出的实践倾向是，尽量让奖励与最终任务完成度一致，同时为关键中间步骤设置轻量 shaping。若中间奖励过强，模型会为了拿分而偏离真实任务；若只有终局奖励，训练又会因为稀疏信号而效率低下。

\begin{importantbox}{奖励设计的工程判据}
有效奖励通常满足三点：可计算、可解释、可扩展。可计算意味着能在流水线上自动生成；可解释意味着失败样本能回溯问题来源；可扩展意味着在新任务上不需要全部重写。
\end{importantbox}

\subsection{数据质量与评估质量的耦合}
讲者在早段就强调 ``data, evaluation and systems'' 同等重要。原因是坏评估会反向污染数据筛选；坏数据又会让评估指标失真，形成负反馈。研发上必须把数据治理和评估治理一起做，而不是分部门串行。

\begin{warningbox}{误区：先把模型训大再补评估}
在 Agent 场景里，这种顺序通常成本更高。模型越大、rollout 越长，后补评估意味着要在更高单次实验成本下重新定位问题，迭代速度会显著下降。
\end{warningbox}

\subsection{本章小结}
Agent 训练阶段的数据单位是轨迹，不是答案；奖励函数要在稀疏终局目标和可学习中间信号间做平衡。最关键的是，数据与评估必须并行设计。


